\documentclass[10pt,a4paper]{article}
\usepackage[margin=23mm]{geometry}
\usepackage{amsmath,amssymb,lmodern}
\setlength{\emergencystretch}{2em}
\title{Short game values are countable\\\large Disproof of Conjecture 00000008873}
\author{ziangni-sys}
\date{4 October 2026}
\begin{document}
\begin{center}
{\Large Short game values are countable}\\[4pt]
{\large Disproof of Conjecture 00000008873}\\[6pt]
ziangni-sys \qquad 4 October 2026
\end{center}
\section*{Claim and actual game class}
Both originals define short games as the inductive class of finite games
and assert that its cardinality is the continuum, realized by binary
game trees. This cardinal assertion is false. We prove that the set of
actual short combinatorial game values has cardinality at most
$\aleph_0$, strictly smaller than $\mathfrak c=2^{\aleph_0}$.

A finite game form consists of two finite lists of previously constructed
game forms: its Left options and Right options. Any finite option set
can be listed, so these descriptions allow arbitrary finite branching.
They are not restricted to one move for each player. Listing options
may introduce order or redundant descriptions; passing to actual game
values identifies these redundancies and cannot increase cardinality.

\section*{Finite descriptions represent every short game}
Use finite syntax with three constructors:
\[
 c ::= \mathsf{nil}\mid\mathsf{cons}(c,c)\mid\mathsf{node}(c,c).
\]
The first two encode finite lists; the last encodes a game with two
option lists. Interpret a node by the actual game construction
$\{L\mid R\}$, using its decoded Left and Right option lists.
For convenience, a non-node used as a game decodes to zero, and a
non-cons used as an option list decodes to the empty list.
All interpretations are short, by induction on the finite syntax.

Conversely, let $x$ be any actual short pre-game. Its Left and Right
move types are finite, and each move leads to another short game.
By induction on shortness, choose finite descriptions of all successor
games. Enumerate each finite move type by
$\{0,\ldots,m-1\}$, collect the chosen descriptions in those finite
lists, and form a node with these two lists. Its interpretation is a
\emph{relabeling} of $x$: the finite enumerations give bijections of
move types, and the induction hypotheses give relabelings of every
corresponding successor. Relabeling implies genuine game equivalence.
Thus every short game value, including all possible option arrangements,
is represented. This is a proved representation result, not an assumed
surjectivity condition.

\section*{Cardinal obstruction}
Finite syntax is countable: encode constructor tags and the two child
codes by a natural-number pairing function, recursively. Injectivity
follows by recovering the tag and child codes. The interpretation map
therefore gives a surjection from a countable type onto all actual short
game values. A surjective image of a countable set is countable, so
\[
 |\mathrm{ShortValues}|\leq\aleph_0<2^{\aleph_0}=\mathfrak c.
\]
The strict inequality is Cantor's theorem. Hence the continuum-cardinality
clause already fails, regardless of the subsequent leaf-depth assertion.
Infinite binary paths can have continuum cardinality, but they are not
finite games in the source's inductive class.

\newpage
\section*{Formal correspondence and completeness}
\raggedright
The Lean project uses Mathlib's actual \texttt{PGame}, its inductive
\texttt{PGame.Short} evidence, \texttt{PGame.Relabelling}, and the actual
\texttt{Game} quotient by game equivalence. It does not replace all
short games by an arbitrarily chosen countable subfamily.

The type \texttt{Code} is the finite syntax above. Mathlib's
\texttt{deriving Countable} handler produces a natural-number encoding
and a kernel-checked injectivity proof. The mutually recursive
\texttt{decode} and \texttt{decodeOptions} interpret codes as actual
pre-games and lists of pre-games. Nodes use \texttt{PGame.ofLists},
whose move types are finite indices of the decoded option lists.
The theorem \texttt{decoded\_short\_and\_options} proves that every
interpretation is short and every decoded option list consists of
short games.

Crucially, \texttt{represents\_every\_short} quantifies over every
actual short \texttt{PGame}, including arbitrary finite move-index
types. By induction on the actual \texttt{Short} evidence, it enumerates
those types with \texttt{Fintype.equivFin}, uses the recursively
represented successors, and constructs an actual \texttt{Relabelling}.
Each move bijection and each successor correspondence is proved.
This removes irrelevant raw index-type labels without assuming that
raw \texttt{PGame} equality has countable many encodings.

The predicate \texttt{IsShortValue v} says that some actual short
pre-game represents the actual game value $v$.
The type \texttt{ShortValue} is the subtype of all such values in the
Mathlib game quotient. The map \texttt{value} sends a code to its
interpretation's actual game value, together with proved shortness.
The theorem \texttt{value\_surjective} proves that every member of this
full subtype is reached, using the representation theorem and
\texttt{Relabelling.equiv}. Its equality is equality in the actual
\texttt{Game} quotient, not a new artificial relation on syntax.

Finally, \texttt{short\_values\_countable} follows from this actual
surjection; \texttt{short\_value\_cardinal\_lt\_continuum} proves the
strict cardinal inequality; and \texttt{conjecture\_8873\_false}
proves that the claimed cardinal equality cannot hold.

\section*{Verification and scope}
The proof allows every finite Left/Right option family; no bound on
branching, game depth, or size is imposed. It covers short game values
under standard Conway equivalence. Finite game forms modulo relabeling
are likewise represented by the same syntax, so distinguishing more
forms than their values does not create continuum many finite games.

Lean 4.19.0 and Mathlib commit
\texttt{c44e0c8ee63ca166450922a373c7409c5d26b00b}
are pinned in the project. Axiom audits of the syntax encoding,
representation theorem, actual value surjection, and final negation
use only \texttt{propext}, \texttt{Classical.choice}, and
\texttt{Quot.sound}. No extra axioms, admitted proofs, or native
evaluation shortcuts are used. No auxiliary computation is required.
\end{document}
